% GENERATED FILE. Edit canonical lessons, then run npm run generate:books.
% canonical-source-hash: fnv1a64:60bd25a93bdc8ed2
% canonical-lessons: PA-C122-bajra, PA-C122-kanak, PA-C122-makki, PA-C122-baigan, PA-C122-gobhi

\chapter{Millet, Wheat, Maize, Aubergine, Cauliflower}
\label{ch:pa-millet-wheat-maize-aubergine-cauliflower}
\hlchaptermodality{122}

\begin{chapteropening}
\textbf{By the end of this chapter:} \emph{I can say millet, wheat, maize, an aubergine and a cauliflower.}

Complete the last lesson of chapter 122: I can say millet, wheat, maize, an aubergine and a cauliflower.
\end{chapteropening}

\section[\texorpdfstring{\pa{ਬਾਜਰਾ} (bājrā)}{ਬਾਜਰਾ (bājrā)}]{\pa{ਬਾਜਰਾ} (bājrā) — millet}

\label{lesson:PA-C122-bajra}



\begin{quote}
Before the new one: say the Punjabi for chickpeas, then the Punjabi for mung beans.
\end{quote}

\subsection*{You'll want to know: \pa{ਬਾਜਰਾ}}
\textbf{\pa{ਬਾਜਰਾ}} — \emph{bājrā} — \textquotedblleft{}millet\textquotedblright{}.

\begin{grammarlens}[title={What you've built}]
One more word for this chapter.
\end{grammarlens}

\subsection*{Your turn}
\begin{itemize}
\raggedright
  \item \emph{Say it:} \emph{bājrā}
  \item \emph{Say it:} \emph{bājrā}, once more
  \item \emph{Say it:} say \emph{bājrā}
  \item \emph{Recall:} say the Punjabi for chickpeas, then the Punjabi for mung beans, then say \emph{bājrā} again
\end{itemize}

\begin{tcolorbox}[breakable,colback=teal!4,colframe=teal!35!black,title={Before you move on}]
What does \pa{ਬਾਜਰਾ} mean? (\textquotedblleft{}Millet\textquotedblright{}.) Say it once more.
\end{tcolorbox}
\section[\texorpdfstring{\pa{ਕਣਕ} (kaṇak)}{ਕਣਕ (kaṇak)}]{\pa{ਕਣਕ} (kaṇak) — wheat}

\label{lesson:PA-C122-kanak}



\begin{quote}
Before the new one: say the Punjabi for mung beans, then the Punjabi for millet.
\end{quote}

\subsection*{You'll want to know: \pa{ਕਣਕ}}
\textbf{\pa{ਕਣਕ}} — \emph{kaṇak} — \textquotedblleft{}wheat\textquotedblright{}.

\begin{grammarlens}[title={What you've built}]
One more word for this chapter.
\end{grammarlens}

\subsection*{Your turn}
\begin{itemize}
\raggedright
  \item \emph{Say it:} \emph{kaṇak}
  \item \emph{Say it:} \emph{kaṇak}, once more
  \item \emph{Say it:} say \emph{kaṇak}
  \item \emph{Recall:} say the Punjabi for mung beans, then the Punjabi for millet, then say \emph{kaṇak} again
\end{itemize}

\begin{tcolorbox}[breakable,colback=teal!4,colframe=teal!35!black,title={Before you move on}]
What does \pa{ਕਣਕ} mean? (\textquotedblleft{}Wheat\textquotedblright{}.) Say it once more.
\end{tcolorbox}
\section[\texorpdfstring{\pa{ਮੱਕੀ} (makkī)}{ਮੱਕੀ (makkī)}]{\pa{ਮੱਕੀ} (makkī) — maize}

\label{lesson:PA-C122-makki}



\begin{quote}
Before the new one: say the Punjabi for millet, then the Punjabi for wheat.
\end{quote}

\subsection*{You'll want to know: \pa{ਮੱਕੀ}}
\textbf{\pa{ਮੱਕੀ}} — \emph{makkī} — \textquotedblleft{}maize\textquotedblright{}.

\begin{grammarlens}[title={What you've built}]
One more word for this chapter.
\end{grammarlens}

\subsection*{Your turn}
\begin{itemize}
\raggedright
  \item \emph{Say it:} \emph{makkī}
  \item \emph{Say it:} \emph{makkī}, once more
  \item \emph{Say it:} say \emph{makkī}
  \item \emph{Recall:} say the Punjabi for millet, then the Punjabi for wheat, then say \emph{makkī} again
\end{itemize}

\begin{tcolorbox}[breakable,colback=teal!4,colframe=teal!35!black,title={Before you move on}]
What does \pa{ਮੱਕੀ} mean? (\textquotedblleft{}Maize\textquotedblright{}.) Say it once more.
\end{tcolorbox}
\section[\texorpdfstring{\pa{ਬੈਂਗਣ} (baĩgaṇ)}{ਬੈਂਗਣ (baĩgaṇ)}]{\pa{ਬੈਂਗਣ} (baĩgaṇ) — an aubergine}

\label{lesson:PA-C122-baigan}



\begin{quote}
Before the new one: say the Punjabi for wheat, then the Punjabi for maize.
\end{quote}

\subsection*{You'll want to know: \pa{ਬੈਂਗਣ}}
\textbf{\pa{ਬੈਂਗਣ}} — \emph{baĩgaṇ} — \textquotedblleft{}an aubergine\textquotedblright{}.

\begin{grammarlens}[title={What you've built}]
One more word for this chapter.
\end{grammarlens}

\subsection*{Your turn}
\begin{itemize}
\raggedright
  \item \emph{Say it:} \emph{baĩgaṇ}
  \item \emph{Say it:} \emph{baĩgaṇ}, once more
  \item \emph{Say it:} say \emph{baĩgaṇ}
  \item \emph{Recall:} say the Punjabi for wheat, then the Punjabi for maize, then say \emph{baĩgaṇ} again
\end{itemize}

\begin{tcolorbox}[breakable,colback=teal!4,colframe=teal!35!black,title={Before you move on}]
What does \pa{ਬੈਂਗਣ} mean? (\textquotedblleft{}An aubergine\textquotedblright{}.) Say it once more.
\end{tcolorbox}
\section[\texorpdfstring{\pa{ਗੋਭੀ} (gobhī)}{ਗੋਭੀ (gobhī)}]{\pa{ਗੋਭੀ} (gobhī) — a cauliflower, a cabbage}

\label{lesson:PA-C122-gobhi}



\begin{quote}
Before the new one: say the Punjabi for maize, then the Punjabi for an aubergine.
\end{quote}

\subsection*{You'll want to know: \pa{ਗੋਭੀ}}
\textbf{\pa{ਗੋਭੀ}} — \emph{gobhī} — \textquotedblleft{}a cauliflower, a cabbage\textquotedblright{}.

\begin{grammarlens}[title={What you've built}]
One more word for this chapter.
\end{grammarlens}

\subsection*{Your turn}
\begin{itemize}
\raggedright
  \item \emph{Say it:} \emph{gobhī}
  \item \emph{Say it:} \emph{gobhī}, once more
  \item \emph{Say it:} say \emph{gobhī}
  \item \emph{Recall:} say the Punjabi for maize, then the Punjabi for an aubergine, then say \emph{gobhī} again
\end{itemize}

\begin{tcolorbox}[breakable,colback=teal!4,colframe=teal!35!black,title={Before you move on}]
What does \pa{ਗੋਭੀ} mean? (\textquotedblleft{}A cauliflower, a cabbage\textquotedblright{}.) Say it once more.
\end{tcolorbox}
